\subsubsection{Observed trade-offs and reproducible evidence}
\label{sec:implemented-generator-results}

Tables~\ref{tab:generator-outcomes}, \ref{tab:generator-tail-outcomes}, and
\ref{tab:generator-speed} summarize
the executed four-family experiment introduced in
Section~\ref{sec:generator-evaluation-protocol}.\footnote{Notebook 6, revision
\texttt{tail\_diagnostics\_v2}. Its complete saved record includes configuration,
seeds, source hashes, reproduction instructions, and the preceding execution.}
All families use 5,000 mechanisms, one channel, one trajectory per mechanism,
and innovation standard deviation 0.05. The original and bounded-domain
families start from uniform histories; Shao--Wu uses certified burn-in;
Heaps uses joint stationary Gaussian histories. The construction rules and
noise distributions are part of each family's definition, so this compares
complete generation policies.

\begin{table}[H]
\centering\small
\begin{tabular}{@{}lrrrr@{}}
\toprule
& \multicolumn{3}{c}{Explosion rate (\%)} & Near-zero band (\%) \\
\cmidrule(lr){2-4}
Generator & $h=96$ & $h=192$ & $h=384$ & $h=384$ \\
\midrule
Original AST & 46.26 & 49.60 & 53.16 & 10.40 \\
Bounded-domain contraction & 0.00 & 0.00 & 0.00 & 44.04 \\
Shao--Wu construction & 0.00 & 0.00 & 0.00 & 23.18 \\
Heaps VAR & 0.00 & 0.00 & 0.00 & 41.54 \\
\bottomrule
\end{tabular}
\caption{Finite-horizon diagnostics, with all 5,000 attempts in each
denominator. Near-zero band means that the final 24 values lie within
$[-0.15,0.15]$ in a non-explosive prefix.}
\label{tab:generator-outcomes}
\end{table}

The rerun reproduces every preceding explosion and near-zero count exactly.
All original-AST failures in this record are finite threshold crossings;
execution stops at their first crossing. Each controlled
family records zero crossings in these 5,000 trials. The notebook's 95\%
Wilson interval has an upper endpoint of approximately 0.077\% for a
zero-event count. The probabilistic interpretation concerns draws from the
configured mechanism/trajectory prior. Analytic domain containment provides
the stronger pathwise statement for the bounded-domain family under its
bounded noise law.

\begin{table}[H]
\centering\small
\begin{tabular}{@{}lrrrrrr@{}}
\toprule
& Sampling + & Initialization & Diagnostics & Total & Complete & Output \\
Generator & certificate (s) & / burn-in (s) & (s) & (s) & paths/s & values/s \\
\midrule
Original AST & 1.65 & 0.43 & 3.35 & 25.16 & 93.10 & 42,014.84 \\
Bounded-domain & 3.58 & 0.59 & 6.63 & 52.52 & 95.19 & 36,554.76 \\
Shao--Wu & 3.48 & 24.19 & 6.15 & 78.75 & 63.49 & 24,381.22 \\
Heaps VAR & 14.19 & 0.42 & 5.85 & 75.84 & 65.93 & 25,316.96 \\
\bottomrule
\end{tabular}
\caption{Recorded serial wall-clock measurements, with BLAS limited to one
thread. Complete paths/s counts trajectories reaching step 384 successfully;
output values/s counts evaluated post-initialization values, including failed
attempts. Both rates divide by total time, including initialization and
the extended diagnostics and their tolerance-sensitivity analysis.}
\label{tab:generator-speed}
\end{table}

These are single-run CPU timings from the recorded Python~3.11.9/NumPy~1.26.4
Windows environment. They characterize the current interpreter and settings;
performance comparisons across machines or parallel backends require repeated,
matched measurements. Heaps is also executed through the common AST evaluator,
which prioritizes a shared execution path over a specialized matrix solver.
Shao--Wu initialization accounts for 1,424,042 burn-in steps across mechanisms,
with median length 303. Its largest recorded retained-history error bound is
$9.999844\times10^{-5}$. The largest Heaps relative Lyapunov residual is
$1.595159\times10^{-15}$.

\begin{table}[H]
\centering\small
\begin{tabular}{@{}lrrrrr@{}}
\toprule
& Persistent & Quiet & Observed & Nonzero & Amplitude \\
Generator & flatness & endpoints & flattening & plateau & decline \\
\midrule
Original AST & 15.70 & 3.84 & 9.18 & 5.02 & 4.94 \\
Bounded-domain & 99.68 & 58.42 & 25.42 & 37.54 & 4.90 \\
Shao--Wu & 75.26 & 68.36 & 0.12 & 47.58 & 0.02 \\
Heaps VAR & 29.36 & 22.98 & 0.22 & 0.02 & 0.30 \\
\bottomrule
\end{tabular}
\caption{Tail diagnostics at $h=384$, as percentages of all 5,000 attempts.
Flatness uses tolerance 0.15 with 24-step windows and 48-step persistence;
flattening and amplitude decline require at least a 50\% reduction from the
initial window, as defined in Section~\ref{sec:generator-evaluation-protocol}.
Labels may overlap.}
\label{tab:generator-tail-outcomes}
\end{table}

Shao--Wu illustrates the distinction between persistent small variation and
observed flattening: 75.26\% satisfy persistent flatness, 68.36\% are quiet
at both observed endpoints, and 0.12\% exhibit the required reduction.
Its 47.58\% of persistent nonzero plateaus also show why the level-independent
diagnostic adds information beyond the near-zero band. For the original AST,
2,342 trajectories reach step 384; persistent flatness is 33.52\% among these
eligible paths, compared with 15.70\% of all attempts.

The tolerance matters substantially. At 0.05, all four families have zero
persistent-flatness events in this sample. At 0.10, their rates are 1.44\%,
67.98\%, 6.00\%, and 1.50\%, respectively, in the table's order.
Thus the 0.15 classification describes variation at the chosen noise-relative
scale. The notebook retains all three tolerance results and Wilson intervals
for the primary diagnostics.


\paragraph{Implications for dataset construction.}
The controlled families substantially reduce observed explosion events while
changing small-amplitude behaviour and generation cost. The added diagnostics
separate observed flattening from quiet initial and terminal windows, and
capture persistent small variation at nonzero levels. Their 0.15 tolerance
is three innovation standard deviations and includes some fluctuating
processes. Changes between finite windows can also arise from stochastic
variation, especially in stationary-initialized families. These measurements
provide operational dataset strata together with their noise scale, tolerance,
and observation duration.

\paragraph{Retained information for the explanation layer.}
The executed equation after transformation is the reference mechanism.
The transformed ASTs, bindings, sampling configurations, seeds, innovation
laws, initialization policies, and certificates make the process auditable.
Several trajectories can share this record. The planned Chapter~5 will use
it to define structural and contribution-based explanation targets over the
forecast horizon. Generator unit tests cover bound propagation, covariance checks,
synchronous updates, seeds, and failure accounting; the notebooks preserve
the empirical evidence and its complete reproduction record.
